% SOURCE_AUTHORITY: sga5_fr_workpass.tex, lines 5894--5931 (LF-normalized unit).
% SOURCE_SHA256: 791F4EFFC5E02832D5D77ED03518C8156D6F07E4C8238B03545DB93D883FBB28
\subsection*{6.8}
Sejam $A$, $B$ umas $\Lambda$-álgebras planas, $X$, $Y$ uns $S$-esquemas, $Z=X\times_SY$, $c:C\to X^2$, $d:D\to Y^2$ umas correspondências,
\[
e=c\times_Sd:E=C\times_SD\longrightarrow Z^2
\]
seu produto. Assim, tem-se
\[
Z^e=X^c\times_SY^d.
\]
Ponhamos $R=A\otimes_{\Lambda}B$. Define-se, como em (III 5.3), o produto tensorial externo
\begin{equation}
\uHom_A(c_1^*L,c_2^!L)\otimes \uHom_B(d_1^*M,d_2^!M)
\xrightarrow{\ \lten_S\ }
\Hom_R(e_1^*(L\lten_SM),e_2^!(L\lten_SM)),
\tag{6.8.1}
\end{equation}
e da comutatividade dos quadrados 5.12.1 e 5.12.3 deduz-se facilmente que, para
\[
u\in\uHom_A(c_1^*L,c_2^!L),\qquad v\in\uHom_B(d_1^*M,d_2^!M),
\]
tem-se, em $H^0(Z^e,KR_{Z^e,L_{\natural}})$,
\begin{equation}
\Tr_R(u\lten_Sv)=\Tr_A(u)\Tr_B(v),
\tag{6.8.2}
\end{equation}
onde o produto do segundo membro é definido pela flecha canônica
\begin{equation}
H^0(X^c,KA_{X^c,L_{\natural}})\otimes
H^0(Y^d,KB_{Y^d,L_{\natural}})
\longrightarrow
H^0(Z^e,KR_{Z^e,L_{\natural}}),
\tag{6.8.3}
\end{equation}
proveniente do isomorfismo
\[
KA_{X^c,L_{\natural}}\lten_S KB_{Y^d,L_{\natural}}\longrightarrow K(A\otimes B)_{Z^e,L_{\natural}}
\]
dado por (III 1.7.6) e a flecha (4) de 5.9.1.
